
\documentclass{article}
\usepackage{colortbl}
\usepackage{makecell}
\usepackage{multirow}
\usepackage{supertabular}

\begin{document}

\newcounter{utterance}

\centering \large Interaction Transcript for game `dond', experiment `semi', episode 18 with Qwen3.6{-}35B{-}A3B{-}FP8{-}without{-}reasoning.
\vspace{24pt}

{ \footnotesize  \setcounter{utterance}{1}
\setlength{\tabcolsep}{0pt}
\begin{supertabular}{c@{$\;$}|p{.15\linewidth}@{}p{.15\linewidth}p{.15\linewidth}p{.15\linewidth}p{.15\linewidth}p{.15\linewidth}}
    \# & $\;$A & \multicolumn{4}{c}{Game Master} & $\;\:$B\\
    \hline

    \theutterance \stepcounter{utterance}  
    & & \multicolumn{4}{p{0.6\linewidth}}{
        \cellcolor[rgb]{0.9,0.9,0.9}{
            \makecell[{{p{\linewidth}}}]{
                \texttt{\tiny{[P1$\langle$GM]}}
                \texttt{Du spiller et kooperativt forhandlingsspil, hvor du skal enes med en anden spiller om, hvordan du deler et sæt genstande mellem jer.} \\
\\ 
\texttt{Regler:} \\
\texttt{(a) Du og den anden spiller får et sæt genstande. Du får også en hemlig værdifunktion, der viser, hvor meget du værdisætter hver type objekt.} \\
\texttt{(b) Du bytter beskeder med den anden spiller for at blive enig om, hvem der får hvad. Du kan sende maksimalt 5 beskeder hver eller afslutte spillet tidligt ved at fremsætte et hemmeligt forslag når som helst.} \\
\texttt{(c) Du skal indsende et hemmeligt forslag, der angiver de genstande, du vil have, formateret i firkantede parenteser som følger: "{[}Forslag: <antal> <objektnavn>, <antal> <objektnavn>, <...>{]}"} \\
\texttt{(d) Hvis dine forslag er komplementære, dvs. der er nok genstande til at opfylde begge forslag, tildeles hver spiller en score baseret på summen af værdierne for de genstande, hun modtog. Ellers får du begge nul point.} \\
\texttt{(e) Dit mål er at maksimere summen af din score og den anden spillers score.} \\
\\ 
\texttt{Lad os starte.} \\
\\ 
\texttt{Sættet af tilgængelige genstande er:} \\
\\ 
\texttt{4 ure, 1 pen, 2 stene} \\
\\ 
\texttt{Din hemmelige værdifunktion er:} \\
\\ 
\texttt{ur: 1, pen: 6, sten: 0} \\
\\ 
\texttt{VIGTIG: Dine beskeder, medmindre det er et hemmeligt forslag, bliver sendt direkte til den anden spiller, så inkluder ikke noget svar på reglerne eller tekst, der annoncerer din besked. Brug det angivne format for at afgive et hemmeligt forslag. Brug ikke firkantede parenteser, når du kommunikerer med den anden spiller, eller det vil blive fortolket som dit hemmelige forslag.} \\
            }
        }
    }
    & \\ \\

    \theutterance \stepcounter{utterance}  
    & \multicolumn{4}{p{0.6\linewidth}}{
        \cellcolor[rgb]{0.9,0.9,0.9}{
            \makecell[{{p{\linewidth}}}]{
                \texttt{\tiny{[P1$\rangle$GM]}}
                \texttt{Hej! Jeg ser, at vi har 4 ure, 1 pen og 2 stene at fordele.} \\
\\ 
\texttt{Baseret på mine værdier (ur: 1, pen: 6, sten: 0) og antagelsen om, at du også har en værdifunktion, vil jeg gerne foreslå en fordeling, der maksimerer den samlede værdi for os begge.} \\
\\ 
\texttt{Da uret har en lav værdi for mig (1), og stenen har ingen værdi (0), vil jeg prioritere at få pennen, hvis den har høj værdi for dig, eller fordele urene ligeligt.} \\
\\ 
\texttt{Lad mig starte med et åbent forslag for at se, hvad du værdsætter højest.} \\
\\ 
\texttt{Hvis du har brug for at kommunikere med mig, så skriv blot din besked her. Jeg er åben for forslag.} \\
\\ 
\texttt{For at teste en potentiel komplementær fordeling, hvor vi begge får noget, kunne vi overveje følgende:} \\
\\ 
\texttt{Jeg foreslår, at jeg tager pennen (hvis den har høj værdi for mig) eller en del af urene. Men da jeg ikke kender din værdifunktion, vil jeg gerne høre, hvad du mener er en fair fordeling.} \\
\\ 
\texttt{Lad os prøve dette forslag:} \\
\\ 
\texttt{Jeg tager 2 ure.} \\
\texttt{Du tager 2 ure, 1 pen og 2 stene.} \\
\\ 
\texttt{Hvad tænker du om det? Eller har du et bedre forslag?} \\
\\ 
\texttt{Hvis du er enig i, at vi skal maksimere den samlede værdi, og du antager, at pennen har høj værdi for dig, kunne det være en god idé, at du tager pennen.} \\
\\ 
\texttt{Lad mig vide, hvad du foretrækker.} \\
            }
        }
    }
    & & \\ \\

    \theutterance \stepcounter{utterance}  
    & & & \multicolumn{2}{p{0.3\linewidth}}{
        \cellcolor[rgb]{0.9,0.9,0.9}{
            \makecell[{{p{\linewidth}}}]{
                \texttt{\tiny{[GM$|$GM]}}
                \texttt{continue} \\
            }
        }
    }
    & & \\ \\

    \theutterance \stepcounter{utterance}  
    & & & \multicolumn{2}{p{0.3\linewidth}}{
        \cellcolor[rgb]{0.9,0.9,0.9}{
            \makecell[{{p{\linewidth}}}]{
                \texttt{\tiny{[GM$|$GM]}}
                \texttt{Hej! Jeg ser, at vi har 4 ure, 1 pen og 2 stene at fordele.} \\
\\ 
\texttt{Baseret på mine værdier (ur: 1, pen: 6, sten: 0) og antagelsen om, at du også har en værdifunktion, vil jeg gerne foreslå en fordeling, der maksimerer den samlede værdi for os begge.} \\
\\ 
\texttt{Da uret har en lav værdi for mig (1), og stenen har ingen værdi (0), vil jeg prioritere at få pennen, hvis den har høj værdi for dig, eller fordele urene ligeligt.} \\
\\ 
\texttt{Lad mig starte med et åbent forslag for at se, hvad du værdsætter højest.} \\
\\ 
\texttt{Hvis du har brug for at kommunikere med mig, så skriv blot din besked her. Jeg er åben for forslag.} \\
\\ 
\texttt{For at teste en potentiel komplementær fordeling, hvor vi begge får noget, kunne vi overveje følgende:} \\
\\ 
\texttt{Jeg foreslår, at jeg tager pennen (hvis den har høj værdi for mig) eller en del af urene. Men da jeg ikke kender din værdifunktion, vil jeg gerne høre, hvad du mener er en fair fordeling.} \\
\\ 
\texttt{Lad os prøve dette forslag:} \\
\\ 
\texttt{Jeg tager 2 ure.} \\
\texttt{Du tager 2 ure, 1 pen og 2 stene.} \\
\\ 
\texttt{Hvad tænker du om det? Eller har du et bedre forslag?} \\
\\ 
\texttt{Hvis du er enig i, at vi skal maksimere den samlede værdi, og du antager, at pennen har høj værdi for dig, kunne det være en god idé, at du tager pennen.} \\
\\ 
\texttt{Lad mig vide, hvad du foretrækker.} \\
            }
        }
    }
    & & \\ \\

    \theutterance \stepcounter{utterance}  
    & & \multicolumn{4}{p{0.6\linewidth}}{
        \cellcolor[rgb]{0.9,0.9,0.9}{
            \makecell[{{p{\linewidth}}}]{
                \texttt{\tiny{[GM$\rangle$P2]}}
                \texttt{Du spiller et kooperativt forhandlingsspil, hvor du skal enes med en anden spiller om, hvordan du deler et sæt genstande mellem jer.} \\
\\ 
\texttt{Regler:} \\
\texttt{(a) Du og den anden spiller får et sæt genstande. Du får også en hemlig værdifunktion, der viser, hvor meget du værdisætter hver type objekt.} \\
\texttt{(b) Du bytter beskeder med den anden spiller for at blive enig om, hvem der får hvad. Du kan sende maksimalt 5 beskeder hver eller afslutte spillet tidligt ved at fremsætte et hemmeligt forslag når som helst.} \\
\texttt{(c) Du skal indsende et hemmeligt forslag, der angiver de genstande, du vil have, formateret i firkantede parenteser som følger: "{[}Forslag: <antal> <objektnavn>, <antal> <objektnavn>, <...>{]}"} \\
\texttt{(d) Hvis dine forslag er komplementære, dvs. der er nok genstande til at opfylde begge forslag, tildeles hver spiller en score baseret på summen af værdierne for de genstande, hun modtog. Ellers får du begge nul point.} \\
\texttt{(e) Dit mål er at maksimere summen af din score og den anden spillers score.} \\
\\ 
\texttt{Lad os starte.} \\
\\ 
\texttt{Sættet af tilgængelige genstande er:} \\
\\ 
\texttt{4 ure, 1 pen, 2 stene} \\
\\ 
\texttt{Din hemmelige værdifunktion er:} \\
\\ 
\texttt{ur: 1, pen: 4, sten: 1} \\
\\ 
\texttt{VIGTIG: Dine beskeder, medmindre det er et hemmeligt forslag, bliver sendt direkte til den anden spiller, så inkluder ikke noget svar på reglerne eller tekst, der annoncerer din besked. Brug det angivne format for at afgive et hemmeligt forslag. Brug ikke firkantede parenteser, når du kommunikerer med den anden spiller, eller det vil blive fortolket som dit hemmelige forslag.} \\
\\ 
\texttt{Hej! Jeg ser, at vi har 4 ure, 1 pen og 2 stene at fordele.} \\
\\ 
\texttt{Baseret på mine værdier (ur: 1, pen: 6, sten: 0) og antagelsen om, at du også har en værdifunktion, vil jeg gerne foreslå en fordeling, der maksimerer den samlede værdi for os begge.} \\
\\ 
\texttt{Da uret har en lav værdi for mig (1), og stenen har ingen værdi (0), vil jeg prioritere at få pennen, hvis den har høj værdi for dig, eller fordele urene ligeligt.} \\
\\ 
\texttt{Lad mig starte med et åbent forslag for at se, hvad du værdsætter højest.} \\
\\ 
\texttt{Hvis du har brug for at kommunikere med mig, så skriv blot din besked her. Jeg er åben for forslag.} \\
\\ 
\texttt{For at teste en potentiel komplementær fordeling, hvor vi begge får noget, kunne vi overveje følgende:} \\
\\ 
\texttt{Jeg foreslår, at jeg tager pennen (hvis den har høj værdi for mig) eller en del af urene. Men da jeg ikke kender din værdifunktion, vil jeg gerne høre, hvad du mener er en fair fordeling.} \\
\\ 
\texttt{Lad os prøve dette forslag:} \\
\\ 
\texttt{Jeg tager 2 ure.} \\
\texttt{Du tager 2 ure, 1 pen og 2 stene.} \\
\\ 
\texttt{Hvad tænker du om det? Eller har du et bedre forslag?} \\
\\ 
\texttt{Hvis du er enig i, at vi skal maksimere den samlede værdi, og du antager, at pennen har høj værdi for dig, kunne det være en god idé, at du tager pennen.} \\
\\ 
\texttt{Lad mig vide, hvad du foretrækker.} \\
            }
        }
    }
    & \\ \\

    \theutterance \stepcounter{utterance}  
    & & & \multicolumn{4}{p{0.6\linewidth}}{
        \cellcolor[rgb]{0.9,0.9,0.9}{
            \makecell[{{p{\linewidth}}}]{
                \texttt{\tiny{[GM$\langle$P2]}}
                \texttt{Hej! Tak for dit forslag.} \\
\\ 
\texttt{Jeg vil gerne være åben og ærlig fra starten for at sikre, at vi finder en fordeling, der maksimerer den samlede værdi for os begge.} \\
\\ 
\texttt{Min hemmelige værdifunktion er:} \\
\texttt{{-} Ur: 1} \\
\texttt{{-} Pen: 4} \\
\texttt{{-} Sten: 1} \\
\\ 
\texttt{Baseret på dette, er pennen den genstand, jeg værdier højest (4 point), mens ure og stene kun giver 1 point hver.} \\
\\ 
\texttt{Hvis jeg tager pennen, får jeg 4 point. De resterende genstande (4 ure og 2 stene) kan du tage. Hvis din værdifunktion er anderledes, vil jeg meget gerne høre, hvad du sætter pris på.} \\
\\ 
\texttt{Mit første hemmelige forslag er:} \\
\texttt{{[}Jeg tager: 1 pen{]}} \\
\texttt{{[}Du tager: 4 ure, 2 stene{]}} \\
\\ 
\texttt{Hvad tænker du om det? Eller har du et andet forslag, du vil byde ind med?} \\
            }
        }
    }
     \\ \\

    \theutterance \stepcounter{utterance}  
    & & & \multicolumn{2}{p{0.3\linewidth}}{
        \cellcolor[rgb]{0.9,0.9,0.9}{
            \makecell[{{p{\linewidth}}}]{
                \texttt{\tiny{[GM$|$GM]}}
                \texttt{proposal does not start with "forslag:"} \\
            }
        }
    }
    & & \\ \\

\end{supertabular}
}

\end{document}
